        \section{估值结论：先判断市场在为什么付钱}
        估值结论不是“电子特气贵所以全部回避”，也不是“钨和氟化工便宜所以全部买入”。当前价格把三类预期混在一起：第一类是钨价、资源管制和库存利润带来的周期重估；第二类是制冷剂配额和价差带来的现金流重估；第三类是WF6/NF3等高纯电子特气通过客户认证后的远期期权。估值工作必须把这三类价值池拆开，否则会把电子特气的稀缺性套给制冷剂，把钨价弹性套给电子气体，或者把主题热度误当成一致预期。

        本章使用独立估值技能生成估值包：每家公司先按业务模型分类，再选择主估值方法和二级校验；随后建立悲观/基本面/乐观三情景、市场预期桥、市场隐含情绪锚、公开券商/公告锚和估值审计。最终目标不是单一公式目标价，而是基本面锚、市场情绪锚和可得外部锚的加权结果。这样处理的目的，是在承认市场已经支付情绪溢价的同时，不让情绪溢价替代订单、利润和客户认证。

        \section{可配置性排序：先看安全边际，再看证据直接性}
        完整估值总表回答的是“今天的股价还允许什么投资动作”。三美股份的最终目标略高于现价，属于回撤配置；翔鹭钨业、新宙邦、永和股份和章源钨业的最终目标低于现价但仍可通过下一季业绩验证，适合中性观察；电子特气和高拥挤标的虽然产业链位置关键，但现价已经显著高于基本面锚，只能放入事件驱动池。排序不是按概念热度，而是按当前价、最终目标、证据质量和证伪速度共同决定。

        \begin{exhibitbox}[完整估值总表]
        \scriptsize
\resizebox{\exhibitwidth}{!}{
\begin{tabular}{llllllllllll}
\toprule
\textbf{代码} & \textbf{名称} & \textbf{方法} & \textbf{现价} & \textbf{市值} & \textbf{EPS26E} & \textbf{悲观} & \textbf{基本面} & \textbf{乐观} & \textbf{最终目标} & \textbf{空间} & \textbf{行动} \\
\midrule
603379 & 三美股份 & PE+配额 & 66.09 & 408.1亿元 & 3.41 & 44.37 & 71.67 & 102.39 & 67.58 & 2.3\% & 回撤配置 \\
002842 & 翔鹭钨业 & PE+周期折扣 & 52.16 & 170.7亿元 & 2.94 & 29.45 & 53.01 & 76.57 & 51.43 & -1.4\% & 中性观察 \\
300037 & 新宙邦 & PE+SOTP & 83.42 & 662.8亿元 & 2.52 & 40.32 & 70.56 & 100.79 & 69.94 & -16.2\% & 中性观察 \\
605020 & 永和股份 & PE+产业链 & 41.22 & 210.7亿元 & 1.47 & 20.57 & 33.79 & 47.02 & 33.67 & -18.3\% & 中性观察 \\
002378 & 章源钨业 & PE+资源储量 & 36.19 & 435.4亿元 & 1.17 & 14.09 & 25.83 & 37.56 & 27.93 & -22.8\% & 中性观察 \\
600160 & 巨化股份 & PE+SOTP & 50.95 & 1583.1亿元 & 1.57 & 22.02 & 34.61 & 47.19 & 36.35 & -28.7\% & 高位风险 \\
002407 & 多氟多 & PE+材料SOTP & 41.51 & 645.9亿元 & 0.97 & 13.51 & 23.17 & 32.82 & 25.47 & -38.7\% & 高位风险 \\
600378 & 昊华科技 & PE+SOTP & 62.30 & 833.6亿元 & 1.00 & 18.02 & 30.03 & 42.04 & 37.08 & -40.5\% & 高位风险 \\
300346 & 南大光电 & PE+材料平台 & 78.12 & 607.6亿元 & 0.69 & 19.44 & 34.71 & 48.60 & 45.44 & -41.8\% & 高位风险 \\
000657 & 中钨高新 & PE+PEG & 99.51 & 2459.0亿元 & 1.49 & 20.87 & 35.78 & 50.69 & 52.66 & -47.1\% & 高位风险 \\
600549 & 厦门钨业 & PE+PB & 84.50 & 2079.9亿元 & 1.67 & 19.98 & 33.31 & 46.63 & 44.52 & -47.3\% & 高位风险 \\
688549 & 中巨芯-U & PS+市场锚 & 32.10 & 487.1亿元 & 0.02 & 5.25 & 10.49 & 16.78 & 16.54 & -48.5\% & 高位风险 \\
688146 & 中船特气 & PE+市场锚 & 385.60 & 2041.4亿元 & 0.83 & 29.11 & 58.23 & 91.50 & 194.27 & -49.6\% & 高位风险 \\
688268 & 华特气体 & PE+平台溢价 & 231.99 & 310.1亿元 & 1.10 & 27.55 & 49.59 & 71.63 & 98.76 & -57.4\% & 高位风险 \\
002971 & 和远气体 & PE+项目弹性 & 66.10 & 163.7亿元 & 0.32 & 5.70 & 11.41 & 17.43 & 24.96 & -62.2\% & 高位风险 \\
603505 & 金石资源 & PE+资源 & 19.01 & 210.9亿元 & 0.24 & 3.36 & 5.52 & 7.68 & 6.90 & -63.7\% & 高位风险 \\
002549 & 凯美特气 & PE+资源化 & 18.99 & 132.1亿元 & 0.06 & 1.16 & 2.19 & 3.21 & 5.75 & -69.7\% & 高位风险 \\
688106 & 金宏气体 & PB/PS+利润恢复 & 34.50 & 192.2亿元 & 0.04 & 0.53 & 0.99 & 1.44 & 7.34 & -78.7\% & 高位风险 \\
\bottomrule
\end{tabular}
}
\sourcenote{AStock估值模型；2026E为Q1可比化估算，非券商一致预期。}
\end{exhibitbox}

表里最重要的不是两位小数目标价，而是行动分层。最终目标低于现价不等于公司没有产业价值，而是当前股价已经提前支付了远期涨价、国产替代、客户认证和流动性溢价。对于钨和制冷剂，下一步要验证的是Q2/Q3利润能否继续承接价格和价差；对于电子特气，下一步要验证的是WF6/NF3等产品是否出现公告级订单、客户结构、价格区间和毛利率改善。没有这些验证，电子特气的高估值应当被定义为期权，而不是低估。

\section{三情景不是装饰：它区分业绩资产和期权资产}
悲观/基本面/乐观三情景的作用，是把当前价格相对基本面锚支付了多少远期期权费显性化。基本面情景代表在现有财报和公开证据下能够解释的价值，乐观情景代表订单、价格、认证和景气度同步兑现，悲观情景代表主题证伪或周期回落。溢价度越高，说明当前价越依赖未来证据；这类标的可以研究，但不能用静态PE简单下结论。

\begin{exhibitbox}[三情景与泡沫度]
\scriptsize
\resizebox{\exhibitwidth}{!}{
\begin{tabular}{lllllll}
\toprule
\textbf{代码} & \textbf{名称} & \textbf{悲观} & \textbf{基本面} & \textbf{乐观} & \textbf{溢价度} & \textbf{二级校验} \\
\midrule
603379 & 三美股份 & 44.37 & 71.67 & 102.39 & -7.8\% & 价差现金流/配额份额/SOTP \\
002842 & 翔鹭钨业 & 29.45 & 53.01 & 76.57 & -1.6\% & PB/资源自给率/钨价敏感性 \\
300037 & 新宙邦 & 40.32 & 70.56 & 100.79 & 18.2\% & SOTP/PS/电子材料收入占比 \\
605020 & 永和股份 & 20.57 & 33.79 & 47.02 & 22.0\% & 价差现金流/配额份额/SOTP \\
002378 & 章源钨业 & 14.09 & 25.83 & 37.56 & 40.1\% & PB/资源自给率/钨价敏感性 \\
600160 & 巨化股份 & 22.02 & 34.61 & 47.19 & 47.2\% & 价差现金流/配额份额/SOTP \\
002407 & 多氟多 & 13.51 & 23.17 & 32.82 & 79.2\% & SOTP/PS/电子材料收入占比 \\
600378 & 昊华科技 & 18.02 & 30.03 & 42.04 & 107.5\% & PS/PB/客户认证与产品收入占比 \\
300346 & 南大光电 & 19.44 & 34.71 & 48.60 & 125.0\% & PS/PB/客户认证与产品收入占比 \\
000657 & 中钨高新 & 20.87 & 35.78 & 50.69 & 178.1\% & PB/资源自给率/钨价敏感性 \\
600549 & 厦门钨业 & 19.98 & 33.31 & 46.63 & 153.7\% & PB/资源自给率/钨价敏感性 \\
688549 & 中巨芯-U & 5.25 & 10.49 & 16.78 & 206.0\% & PSG/客户认证/现金消耗 \\
688146 & 中船特气 & 29.11 & 58.23 & 91.50 & 562.2\% & PS/PB/客户认证与产品收入占比 \\
688268 & 华特气体 & 27.55 & 49.59 & 71.63 & 367.8\% & PS/PB/客户认证与产品收入占比 \\
002971 & 和远气体 & 5.70 & 11.41 & 17.43 & 479.4\% & PS/PB/客户认证与产品收入占比 \\
603505 & 金石资源 & 3.36 & 5.52 & 7.68 & 244.3\% & SOTP/PS/电子材料收入占比 \\
002549 & 凯美特气 & 1.16 & 2.19 & 3.21 & 768.8\% & PS/PB/客户认证与产品收入占比 \\
688106 & 金宏气体 & 0.53 & 0.99 & 1.44 & 3391.4\% & PS/PB/客户认证与产品收入占比 \\
\bottomrule
\end{tabular}
}
\end{exhibitbox}

这张表把板块分成两类。三美股份、翔鹭钨业这类标的的当前价接近基本面锚，主要矛盾是等待更好的买点或财报确认；中船特气、中巨芯-U、华特气体、金宏气体、凯美特气等标的的溢价度显著高，主要矛盾不是“产业逻辑有没有”，而是市场锚能否被订单、客户和收入占比证明。若只有题材传播而没有公告和财报，估值会从市场锚向基本面锚回落。

\section{市场预期桥：把题材翻译成财报门槛}
市场预期桥回答“投资者到底在为什么付钱”。钨链支付的是资源稀缺和战略管制，要求钨价维持高位并进入矿端或钨粉利润；制冷剂链支付的是配额和价差现金流，要求R32/R125/R134a景气不快速回落；电子特气支付的是高纯品类、晶圆厂认证和进口替代持续时间，要求收入、毛利率和客户验证同步出现。只要预期无法落到这些财报项目，市场锚就应下调。

\begin{exhibitbox}[市场预期估值桥]
\scriptsize
\resizebox{\exhibitwidth}{!}{
\begin{tabular}{lllllllll}
\toprule
\textbf{代码} & \textbf{名称} & \textbf{现价} & \textbf{2026E收入} & \textbf{收入增速} & \textbf{归母/EPS} & \textbf{倍数} & \textbf{基本面锚} & \textbf{驱动} \\
\midrule
603379 & 三美股份 & 66.09 & 58.2亿元 & -0.4\% & 21.1亿元/3.41 & 21.0x & 71.67 & 制冷剂价差+配额现金流 \\
002842 & 翔鹭钨业 & 52.16 & 44.2亿元 & 83.3\% & 9.6亿元/2.94 & 18.0x & 53.01 & 价格高位+资源稀缺+战略管制 \\
300037 & 新宙邦 & 83.42 & 140.0亿元 & 45.3\% & 20.0亿元/2.52 & 28.0x & 70.56 & 材料SOTP+电子级产品验证 \\
605020 & 永和股份 & 41.22 & 55.2亿元 & 6.0\% & 7.5亿元/1.47 & 23.0x & 33.79 & 制冷剂价差+配额现金流 \\
002378 & 章源钨业 & 36.19 & 97.4亿元 & 87.3\% & 14.1亿元/1.17 & 22.0x & 25.83 & 价格高位+资源稀缺+战略管制 \\
600160 & 巨化股份 & 50.95 & 250.7亿元 & -7.1\% & 48.9亿元/1.57 & 22.0x & 34.61 & 制冷剂价差+配额现金流 \\
002407 & 多氟多 & 41.51 & 128.6亿元 & 36.3\% & 15.0亿元/0.97 & 24.0x & 23.17 & 材料SOTP+电子级产品验证 \\
600378 & 昊华科技 & 62.30 & 184.0亿元 & 10.2\% & 13.4亿元/1.00 & 30.0x & 30.03 & 客户认证+进口替代+高纯电子特气涨价 \\
300346 & 南大光电 & 78.12 & 28.8亿元 & 11.3\% & 5.4亿元/0.69 & 50.0x & 34.71 & 客户认证+进口替代+高纯电子特气涨价 \\
000657 & 中钨高新 & 99.51 & 280.3亿元 & 58.9\% & 36.8亿元/1.49 & 24.0x & 35.78 & 价格高位+资源稀缺+战略管制 \\
600549 & 厦门钨业 & 84.50 & 583.1亿元 & 26.0\% & 41.0亿元/1.67 & 20.0x & 33.31 & 价格高位+资源稀缺+战略管制 \\
688549 & 中巨芯-U & 32.10 & 15.9亿元 & 31.4\% & 0.3亿元/0.02 & 10.0x & 10.49 & 客户认证+进口替代+高纯电子特气涨价 \\
688146 & 中船特气 & 385.60 & 30.5亿元 & 34.9\% & 4.4亿元/0.83 & 70.0x & 58.23 & 客户认证+进口替代+高纯电子特气涨价 \\
688268 & 华特气体 & 231.99 & 16.7亿元 & 17.7\% & 1.5亿元/1.10 & 45.0x & 49.59 & 客户认证+进口替代+高纯电子特气涨价 \\
002971 & 和远气体 & 66.10 & 17.0亿元 & 2.4\% & 0.8亿元/0.32 & 36.0x & 11.41 & 客户认证+进口替代+高纯电子特气涨价 \\
603505 & 金石资源 & 19.01 & 40.0亿元 & 2.9\% & 2.7亿元/0.24 & 23.0x & 5.52 & 材料SOTP+电子级产品验证 \\
002549 & 凯美特气 & 18.99 & 6.4亿元 & 2.2\% & 0.4亿元/0.06 & 34.0x & 2.19 & 客户认证+进口替代+高纯电子特气涨价 \\
688106 & 金宏气体 & 34.50 & 28.8亿元 & 3.8\% & 0.2亿元/0.04 & 26.0x & 0.99 & 客户认证+进口替代+高纯电子特气涨价 \\
\bottomrule
\end{tabular}
}
\end{exhibitbox}

这张桥也解释了为什么不能把全板块统一用PE比较。中巨芯-U接近零EPS，使用PS和现金消耗校验比PE更合理；中船特气和华特气体有电子特气平台价值，但当前价隐含的增长持续时间很长；巨化、三美、永和的主要利润仍来自制冷剂，电子材料只能作为SOTP期权；新宙邦需要同时拆分有机氟、电子氟化液和锂电材料周期。方法匹配业务模型，比选一个看似统一的行业倍数更重要。

\section{情绪锚：承认市场共识，但不让它绑架目标价}
强主题行情里，市场锚有信息含量：成交额、价格趋势、稀缺叙事和公告催化会让股票长期高于静态基本面锚。模型保留市场锚，是为了把这种现实纳入最终目标；但市场锚只解释价格，不证明价值。基本面权重越低，后续证伪速度越快；如果公告反复提示订单有限、收入占比低或无扩产计划，市场权重必须下调。

\begin{exhibitbox}[市场隐含情绪锚]
\scriptsize
\resizebox{\exhibitwidth}{!}{
\begin{tabular}{llllllllll}
\toprule
\textbf{代码} & \textbf{名称} & \textbf{内在值} & \textbf{隐含倍数} & \textbf{交易} & \textbf{情绪} & \textbf{市场锚} & \textbf{权重F/M/S} & \textbf{最终目标} & \textbf{溢价} \\
\midrule
603379 & 三美股份 & 71.67 & 19.4x & 流动性一般 & 基本面主导 & 51.22 & 80\%/20\%/0\% & 67.58 & -7.8\% \\
002842 & 翔鹭钨业 & 53.01 & 17.7x & 中高成交额 & 主题支持但需业绩 & 46.68 & 75\%/25\%/0\% & 51.43 & -1.6\% \\
300037 & 新宙邦 & 70.56 & 33.1x & 中高成交额 & 基本面主导 & 67.15 & 82\%/18\%/0\% & 69.94 & 18.2\% \\
605020 & 永和股份 & 33.79 & 28.1x & 中高成交额 & 基本面主导 & 33.18 & 80\%/20\%/0\% & 33.67 & 22.0\% \\
002378 & 章源钨业 & 25.83 & 30.8x & 中高成交额 & 主题支持但需业绩 & 32.39 & 68\%/32\%/0\% & 27.93 & 40.1\% \\
600160 & 巨化股份 & 34.61 & 32.4x & 高成交额/拥挤 & 基本面主导 & 42.54 & 78\%/22\%/0\% & 36.35 & 47.2\% \\
002407 & 多氟多 & 23.17 & 43.0x & 高成交额/拥挤 & 基本面主导 & 34.66 & 80\%/20\%/0\% & 25.47 & 79.2\% \\
600378 & 昊华科技 & 30.03 & 62.2x & 高成交额/拥挤 & 强情绪溢价 & 58.25 & 75\%/25\%/0\% & 37.08 & 107.5\% \\
300346 & 南大光电 & 34.71 & 112.5x & 高成交额/拥挤 & 强情绪溢价 & 73.04 & 72\%/28\%/0\% & 45.44 & 125.0\% \\
000657 & 中钨高新 & 35.78 & 66.7x & 高成交额/拥挤 & 强情绪溢价 & 92.05 & 70\%/30\%/0\% & 52.66 & 178.1\% \\
600549 & 厦门钨业 & 33.31 & 50.7x & 高成交额/拥挤 & 强情绪溢价 & 78.16 & 75\%/25\%/0\% & 44.52 & 153.7\% \\
688549 & 中巨芯-U & 10.49 & 30.6x & 中高成交额 & 强情绪溢价 & 27.77 & 65\%/35\%/0\% & 16.54 & 206.0\% \\
688146 & 中船特气 & 58.23 & 463.6x & 高成交额/拥挤 & 强情绪溢价 & 360.54 & 55\%/45\%/0\% & 194.27 & 562.2\% \\
688268 & 华特气体 & 49.59 & 210.5x & 中高成交额 & 强情绪溢价 & 209.95 & 60\%/30\%/10\% & 98.76 & 367.8\% \\
002971 & 和远气体 & 11.41 & 208.6x & 中高成交额 & 强情绪溢价 & 59.82 & 72\%/28\%/0\% & 24.96 & 479.4\% \\
603505 & 金石资源 & 5.52 & 79.2x & 流动性一般 & 强情绪溢价 & 14.73 & 85\%/15\%/0\% & 6.90 & 244.3\% \\
002549 & 凯美特气 & 2.19 & 295.4x & 中高成交额 & 强情绪溢价 & 16.43 & 75\%/25\%/0\% & 5.75 & 768.8\% \\
688106 & 金宏气体 & 0.99 & 907.8x & 中高成交额 & 强情绪溢价 & 29.84 & 78\%/22\%/0\% & 7.34 & 3391.4\% \\
\bottomrule
\end{tabular}
}
\end{exhibitbox}

对电子特气，高市场锚意味着“持有研究权”，不是“自动买入权”。中船特气可以因为WF6/NF3直接性和高纯认证获得较高市场权重，但公告若不能继续提供订单、价格和客户验证，最终目标应向基本面锚收敛；中巨芯-U拥有直接产能证据，但近零EPS阶段必须用PS和现金流消耗观察；金宏气体、凯美特气和和远气体属于扩散线，除非产品收入和客户证据变清楚，否则市场锚权重不能继续提高。

\section{券商和公告锚：缺失本身就是风险信息}
公开券商目标价和一致预期并不完整，不能把缺失字段用主观假设填满。能引用的公开锚主要包括华特气体历史PDF目标价、制冷剂行业配额点评，以及多家公司关于WF6产能、订单、扩产和收入占比的公告或媒体转述。券商锚缺失不是空白页，而是证据等级降低；在这种情况下，AStock目标价必须更依赖财报可验证项和公告边界。

\begin{exhibitbox}[公开券商/公告锚对比]
\scriptsize
\resizebox{\exhibitwidth}{!}{
\begin{tabular}{llllllll}
\toprule
\textbf{代码} & \textbf{名称} & \textbf{来源} & \textbf{日期} & \textbf{评级/目标} & \textbf{外部目标} & \textbf{方法/假设} & \textbf{质量} \\
\midrule
603379 & 三美股份 & 国信证券配额点评转载 & 2025-12-11 & 行业看好，单股目标未披露 & 未披露 & 配额景气与价差框架 & 高 \\
002842 & 翔鹭钨业 & 未披露 & 未披露 & 未披露 & 未披露 & 未披露 & 中 \\
300037 & 新宙邦 & 国信证券配额点评转载 & 2025-12-11 & 行业看好，单股目标未披露 & 未披露 & 配额景气与价差框架 & 中 \\
605020 & 永和股份 & 国信证券配额点评转载 & 2025-12-11 & 行业看好，单股目标未披露 & 未披露 & 配额景气与价差框架 & 中 \\
002378 & 章源钨业 & 未披露 & 未披露 & 未披露 & 未披露 & 未披露 & 高 \\
600160 & 巨化股份 & 国信证券配额点评转载 & 2025-12-11 & 行业看好，单股目标未披露 & 未披露 & 配额景气与价差框架 & 高 \\
002407 & 多氟多 & 国信证券配额点评转载 & 2025-12-11 & 行业看好，单股目标未披露 & 未披露 & 配额景气与价差框架 & 中 \\
600378 & 昊华科技 & 公告/媒体回应 & 2026-06 & 无目标价，公告边界证据 & 未披露 & 产能/订单/收入占比边界 & 高 \\
300346 & 南大光电 & 公告/媒体回应 & 2026-06 & 无目标价，公告边界证据 & 未披露 & 产能/订单/收入占比边界 & 中 \\
000657 & 中钨高新 & 未披露 & 未披露 & 未披露 & 未披露 & 未披露 & 中 \\
600549 & 厦门钨业 & 未披露 & 未披露 & 未披露 & 未披露 & 未披露 & 高 \\
688549 & 中巨芯-U & 公告/媒体回应 & 2026-06 & 无目标价，公告边界证据 & 未披露 & 产能/订单/收入占比边界 & 高 \\
688146 & 中船特气 & 公告/媒体回应 & 2026-06 & 无目标价，公告边界证据 & 未披露 & 产能/订单/收入占比边界 & 高 \\
688268 & 华特气体 & 华泰证券PDF & 2024-11-07 & 买入/目标60.21元 & 60.21 & PE，2026E EPS 2.85元 & 高 \\
002971 & 和远气体 & 公告/媒体回应 & 2026-06 & 无目标价，公告边界证据 & 未披露 & 产能/订单/收入占比边界 & 中 \\
603505 & 金石资源 & 国信证券配额点评转载 & 2025-12-11 & 行业看好，单股目标未披露 & 未披露 & 配额景气与价差框架 & 中 \\
002549 & 凯美特气 & 公告/媒体回应 & 2026-06 & 无目标价，公告边界证据 & 未披露 & 产能/订单/收入占比边界 & 中 \\
688106 & 金宏气体 & 公告/媒体回应 & 2026-06 & 无目标价，公告边界证据 & 未披露 & 产能/订单/收入占比边界 & 中 \\
\bottomrule
\end{tabular}
}
\end{exhibitbox}

因此，外部目标价在本章里只是校验锚，不是替代估值。华特气体历史目标价明显低于当前价格，说明市场已经重新支付电子特气和国产替代溢价；氟化工公司只有行业景气锚，没有完整单股目标价；中船特气、中巨芯和昊华科技的公告边界更像风险锚，提醒不要把WF6价格传闻无差别转成利润预测。

\section{方法和失效条件：每个标的都要知道自己怎么被证伪}
方法桥把每家公司放回自己的商业模式。钨链不能只看半导体概念，要看资源自给率、钨价敏感性和PB校验；制冷剂公司要看配额、价差和现金流；材料SOTP公司要拆出电子材料收入占比；电子特气平台要看客户认证、产品组合、收入占比和PS/PB校验。估值失效条件不是附录文字，而是后续调研和调仓的触发器。

\begin{exhibitbox}[方法与失效条件桥]
\scriptsize
\resizebox{\exhibitwidth}{!}{
\begin{tabular}{llllll}
\toprule
\textbf{代码} & \textbf{名称} & \textbf{业务模型} & \textbf{主方法} & \textbf{二级校验} & \textbf{失效条件} \\
\midrule
603379 & 三美股份 & 配额现金流+材料期权 & PE+配额 & 价差现金流/配额份额/SOTP & 制冷剂价格低于配额景气预期，或含氟材料/电子级产品认证低于预期 \\
002842 & 翔鹭钨业 & 资源/周期盈利 & PE+周期折扣 & PB/资源自给率/钨价敏感性 & 钨价快速回落、矿端供给释放或下游硬质合金需求承压导致Q2利润环比下滑 \\
300037 & 新宙邦 & 化工材料SOTP & PE+SOTP & SOTP/PS/电子材料收入占比 & 制冷剂价格低于配额景气预期，或含氟材料/电子级产品认证低于预期 \\
605020 & 永和股份 & 配额现金流+材料期权 & PE+产业链 & 价差现金流/配额份额/SOTP & 制冷剂价格低于配额景气预期，或含氟材料/电子级产品认证低于预期 \\
002378 & 章源钨业 & 资源/周期盈利 & PE+资源储量 & PB/资源自给率/钨价敏感性 & 钨价快速回落、矿端供给释放或下游硬质合金需求承压导致Q2利润环比下滑 \\
600160 & 巨化股份 & 配额现金流+材料期权 & PE+SOTP & 价差现金流/配额份额/SOTP & 制冷剂价格低于配额景气预期，或含氟材料/电子级产品认证低于预期 \\
002407 & 多氟多 & 化工材料SOTP & PE+材料SOTP & SOTP/PS/电子材料收入占比 & 制冷剂价格低于配额景气预期，或含氟材料/电子级产品认证低于预期 \\
600378 & 昊华科技 & 认证驱动电子材料平台 & PE+SOTP & PS/PB/客户认证与产品收入占比 & WF6价格/订单缺乏公告验证，或单一产品预期被证伪导致估值锚下移 \\
300346 & 南大光电 & 认证驱动电子材料平台 & PE+材料平台 & PS/PB/客户认证与产品收入占比 & WF6价格/订单缺乏公告验证，或单一产品预期被证伪导致估值锚下移 \\
000657 & 中钨高新 & 资源/周期盈利 & PE+PEG & PB/资源自给率/钨价敏感性 & 钨价快速回落、矿端供给释放或下游硬质合金需求承压导致Q2利润环比下滑 \\
600549 & 厦门钨业 & 资源/周期盈利 & PE+PB & PB/资源自给率/钨价敏感性 & 钨价快速回落、矿端供给释放或下游硬质合金需求承压导致Q2利润环比下滑 \\
688549 & 中巨芯-U & 收入高增/近零EPS材料平台 & PS+市场锚 & PSG/客户认证/现金消耗 & WF6价格/订单缺乏公告验证，或单一产品预期被证伪导致估值锚下移 \\
688146 & 中船特气 & 认证驱动电子材料平台 & PE+市场锚 & PS/PB/客户认证与产品收入占比 & WF6价格/订单缺乏公告验证，或单一产品预期被证伪导致估值锚下移 \\
688268 & 华特气体 & 认证驱动电子材料平台 & PE+平台溢价 & PS/PB/客户认证与产品收入占比 & WF6价格/订单缺乏公告验证，或单一产品预期被证伪导致估值锚下移 \\
002971 & 和远气体 & 认证驱动电子材料平台 & PE+项目弹性 & PS/PB/客户认证与产品收入占比 & WF6价格/订单缺乏公告验证，或单一产品预期被证伪导致估值锚下移 \\
603505 & 金石资源 & 化工材料SOTP & PE+资源 & SOTP/PS/电子材料收入占比 & 制冷剂价格低于配额景气预期，或含氟材料/电子级产品认证低于预期 \\
002549 & 凯美特气 & 认证驱动电子材料平台 & PE+资源化 & PS/PB/客户认证与产品收入占比 & WF6价格/订单缺乏公告验证，或单一产品预期被证伪导致估值锚下移 \\
688106 & 金宏气体 & 认证驱动电子材料平台 & PB/PS+利润恢复 & PS/PB/客户认证与产品收入占比 & WF6价格/订单缺乏公告验证，或单一产品预期被证伪导致估值锚下移 \\
\bottomrule
\end{tabular}
}
\end{exhibitbox}

\section{最终行动框架}
回撤配置池只放基本面锚相对扎实、价格没有严重透支的标的。三美股份的制冷剂现金流和配额逻辑最接近这一类，但买点仍取决于价差持续性和回撤后的安全边际；翔鹭钨业虽然为中性观察，但如果钨价和Q2利润继续兑现，估值可以重新靠近配置池。

中性观察池关注基本面锚能否被财报抬高。新宙邦、永和股份和章源钨业都不是简单看多或看空：新宙邦需要电子氟化液和有机氟材料给出SOTP贡献，永和股份需要制冷剂价差和一体化兑现，章源钨业需要钨价高位转化为利润和现金流。Q2结果如果低于门槛，最终目标应下调；如果利润和价格信号超预期，基本面锚可以上修。

事件验证池用于跟踪高弹性，不用于无条件配置。中船特气、中巨芯-U、华特气体、南大光电、金宏气体、凯美特气和和远气体的共同问题，是市场已经提前支付电子特气国产替代期权。只有当订单、客户、产品收入占比、价格和毛利率逐步被公告或财报验证时，市场锚才有资格保留；否则这些股票的研究结论应从“稀缺资产”下调为“高位情绪溢价等待证伪”。
